\subsubsection{Translation of Binary Boolean Operators}
Integer-C conjunction and disjunction evaluate their left operand first and may skip their right operand.
This matters when the right operand calls a function or would raise an arithmetic error.
The F\# backend implements this source-language behavior with JVM conditional branches.

Let $(L_1,s_1)$ and $(L_2,s_2)$ be the instruction lists and stack bounds for the two operands.
For conjunction, emit the following instructions, with fresh labels $\ell_0$ and $\ell_e$:
\[
\begin{aligned}
&L_1;\ \mytt{ifeq}\ \ell_0;\ L_2;\ \mytt{ifeq}\ \ell_0;\\
&\mytt{iconst\_1};\ \mytt{goto}\ \ell_e;\
\ell_0\mytt{:};\ \mytt{iconst\_0};\ \ell_e\mytt{:}.
\end{aligned}
\]
Each \mytt{ifeq} consumes the tested operand. A zero value branches immediately to the false result;
the second operand is evaluated only after a nonzero first value.
The two result paths leave exactly one normalized Boolean, zero or one, on the stack.
Since the left value has been consumed before the right operand runs, the maximum extra stack height is
\[
\max(1,s_1,s_2).
\]
For disjunction, use \mytt{ifne} for both tests, emit \mytt{iconst\_0} after the second test,
and emit \mytt{iconst\_1} at the decisive label. Its stack bound is the same.
The complete \texttt{compileExpr} implementation below uses precisely these two patterns.

The following retained derivation is a comparison with an eager implementation for operands already
normalized to zero or one. Eager evaluation is suitable only when evaluating the right operand is intended;
it does not define Integer-C's short-circuit operators. Its larger stack bound follows because the left
value remains on the stack during the evaluation of the right operand.
